\section{Cube conventions and elementary reductions}\label{sec:preliminaries}
All averages and inner products below use the uniform measure on the relevant cube. For a real function \(u\) on \(\{-1,1\}^n\), write
\[
\chi_S(x)=\prod_{i\in S}x_i,\qquad
\widehat u(S)=\E[u\chi_S],\qquad
u=\sum_{S\subseteq[n]}\widehat u(S)\chi_S.
\]
The functions \(\chi_S\) form an orthonormal basis: independence makes
\(\E\chi_S\chi_T\) equal to \(1\) if \(S=T\) and to \(0\) otherwise.
Consequently,
\[
\Var(u)=\sum_{S\ne\varnothing}\widehat u(S)^2.
\]
For later use set \(W_k(u)=\sum_{|S|=k}\widehat u(S)^2\).

Let \(\E_i\) average only coordinate \(i\), and define
\[
D_i=I-\E_i,\qquad \mathcal N=\sum_{i=1}^n D_i.
\]
Thus \(D_i u(x)=(u(x)-u(x^{(i)}))/2\), where \(x^{(i)}\) reverses coordinate \(i\).
The projections \(D_i\) are self-adjoint and satisfy
\begin{equation}\label{eq:dirichlet}
\E[u\mathcal N v]
=\sum_i\E[D_i u\,D_i v]
=\sum_S |S|\,\widehat u(S)\widehat v(S).
\end{equation}
Independence in the noise definition gives
\(T_\rho\chi_S=\rho^{|S|}\chi_S\). Hence, for \(\tau\ge0\),
\(T_{e^{-\tau}}=e^{-\tau\mathcal N}\).
These elementary facts are also standard in the analysis of Boolean functions~\cite{ODonnell14}; their normalizations are specified here to fix every factor of two.

We first remove the endpoints and symmetries.
The identity
\[
T_{-\rho}f(x)=T_\rho f(-x)
\]
shows that it suffices to consider \(\rho\ge0\). At \(\rho=0\), the left side of Equation~\eqref{eq:main} is zero. At \(\rho=1\), it equals \(H(m)\le1\). For a constant \(f\), it is again zero.
For \(0<\rho<1\) and nonconstant \(f\), every input has positive noise probability, so
\[
-1<T_\rho f(x)<1\qquad\text{for every }x.
\]
All differentiations involving \(H\) below are therefore legitimate. Finally,
replacing \(f\) by \(-f\) preserves both sides of Equation~\eqref{eq:main} and reverses \(m\).

